\paragraph{\trIndex{}理论与\trLagrangianGrassmannian。}\label{sub:index-theory-grassmannian}
设 $V=(V,\varphi_V)$ 是一个带\trSymplecticStructure{}的 $2n$ 维\trVectorSpace，而
$Gr(V)\cong\mathrm{U}_n/\mathrm{O}_n$ 是它的由（无\trOrientation）\trLinear\trLagrange{}子空间构成
的\trGrassmann\trManifold。一个映射 $B\to Gr(V)$ 无非是平凡辛\trVectorBundle{}
$B\times V\to B$ 的一个\trLagrange\trSubbundle，因此 $Gr(V)$ 的\trCohomology{}类可以被
看作\textbf{相对示性类}（relative characteristic classes）。我们关心的是\trMaslovClass{}
$\mu\in H^1(Gr(V))$，以及第二\trStiefelWhitney{}类 $w_2\in H^2(Gr(V);\ZZ/2)$——后者由
全体无\trOrientation{}子空间的寻常\trGrassmann\trManifold{}诱导而来（第一%
\trStiefelWhitney{}类是 $\mu$ 的 mod 2 约化，因而是多余的）。与 $\mu$ 的\trPairing{}是
一个\trIsomorphism{} $\pi_1(Gr(V))\to\ZZ$。当 $n\ge 3$ 时，$w_2$ 类似地给出一个%
\trIsomorphism{} $\pi_2(Gr(V))\to\ZZ/2$。例外的低维情形是：$n=1$ 时
$\mathrm{U}_1/\mathrm{O}_1=S^1$，故 $\pi_2(Gr(V))=0$；$n=2$ 时
$\mathrm{U}_2/\mathrm{O}_2=S^1\times_{\ZZ/2}S^2$，故 $\pi_2(Gr(V))\cong\ZZ$。最后提到的
这个群由\trDoubleCover{} $\mathrm{U}_2/\mathrm{SO}_2$ 上的\trEuler{}类所检测；由于后者在
模 2 下约化为第二\trStiefelWhitney{}类，$w_2\colon\pi_2(\mathrm{U}_2/\mathrm{O}_2)\to\ZZ/2$
仍是满的。
我们还需要考察带基与自由的\trLoopSpace{} $\Omega\,Gr(V)$、$\scr{L}\,Gr(V)$。它们有
$\ZZ$ 那么多的\trConnectedComponent{} $\Omega_kGr(V)$ 或 $\scr{L}_kGr(V)$，由%
\trMaslovIndex{} $\mu(\rho)=\langle\mu,[\rho]\rangle=k$ 相区分。当 $n\ge 3$ 时容易看出：
$\pi_1(\Omega_kGr(V))\cong\ZZ/2$，而
$\pi_1(\scr{L}_kGr(V))\cong\pi_2(Gr(V))\times\pi_1(Gr(V))\cong\ZZ/2\times\ZZ$。
特别地，
\begin{equation}\label{eq:free-loop-cohomology}
H^1(\scr{L}_kGr(V);\ZZ/2)\cong\ZZ/2\oplus\ZZ/2
\end{equation}
带有显式\trGenerator{} $T(w_2)$ 与（$U(\mu)$ 的 mod 2 约化）。这里的记号约定是：
$(T,U)$ 是\trPullback{}映射 $H^*(Gr(V))\to H^*(S^1\times\scr{L}\,Gr(V))\cong
H^{*-1}(\scr{L}\,Gr(V))\oplus H^*(\scr{L}\,Gr(V))$ 的两个分量。
为了与\trIndex{}理论建立联系，做法如下。在 $V$ 上取定一个 $\varphi_V$-\trCompatible{}的%
\trComplexStructure{} $I_V$、一个恰有一个\trBoundaryComponent{}的\trCompact\trRiemann\trSurface，以及
一个\trOrientation{}的等同 $\partial S\cong S^1$。对任意 $\rho\in\scr{L}\,Gr(V)$，令
$D_{S,\rho}$ 为作用在平凡复\trVectorBundle{} $E=S\times V$ 的诸\trSection{}上的
$\bar\partial$-算子，其\trBoundaryCondition{}由\trLagrange\trSubbundle{} $F_\rho\subset E|\partial S$
给出：$(F_\rho)_{\mathrm{exp}(2\pi is)}=\rho(s)$。让 $\rho$ 变动而考察算子族 $D_{S,\rho}$，
便得到一个\trDeterminantLine\trLineBundle{} $\underline{\det}_S\to\scr{L}\,Gr(V)$。
\begin{lemma}\label{thm:det-line-bundle-loop-space}
$\mathrm{index}(D_{S,\rho})=n\chi(S)+\mu(\rho)$。此外，\trDeterminantLine\trLineBundle{}的%
\trIsomorphism{}类由
\begin{equation}\label{eq:det-bundle-w1}
w_1(\underline{\det}_S)=T(w_2)+(T(\mu)-1)\,U(\mu)
\end{equation}
给出。
\end{lemma}
\begin{proof}
第一个陈述是众所周知的带边\trSurface{}的\trRiemann--\trRoch{}定理。第二个陈述同样可以
被看作“Real”算子族的\trIndex{}定理的一个实例 [9]。此外，还有一个不那么直接但更
显式的论证，基于如下几件素材：
(i) \emph{\trPinch{}掉边界}。假设 $S$ 的\trGenus{} $>0$。此时可以把它看作一个\trDisk{} $D$ 与
一个同\trGenus{}的闭\trRiemann\trSurface{}的\trConnectedSum。闭\trSurface{}上的 $\bar\partial$-算子
显然是 $\CC$-\trLinear{}的，因而其\trDeterminantLine{}有一个\trCanonical{}\trTrivialization。由
\eqref{eq:interior-gluing} 便得到一个\trIsomorphism{}
$\underline{\det}_S\cong\underline{\det}_D$。
(ii) \emph{\trDirectSum{}公式}。设
$\Delta\colon Gr(V)\times Gr(W)\to Gr(V\oplus W)$ 为\trDirectSum{}运算
$(\Lambda_1,\Lambda_2)\mapsto\Lambda_1\oplus\Lambda_2$，并令
$\scr{L}_{k,l}\Delta\colon\scr{L}_kGr(V)\times\scr{L}_lGr(W)\to
\scr{L}_{k+l}Gr(V\oplus W)$ 为自由\trLoopSpace{}上被诱导的映射。由
\eqref{eq:det-direct-sum}，
\begin{equation*}
(\scr{L}_{k,l}\Delta)^*\,\underline{\det}_{S,V\oplus W}
\cong\underline{\det}_{S,V}\boxtimes\underline{\det}_{S,W}
\end{equation*}
（这里我们在记号中补上了所涉的\trVectorSpace，以澄清关系）。另一方面，用\trWhitney{}和
公式，
\begin{equation*}
\begin{aligned}
(\scr{L}_{k,l}\Delta)^*(T(w_2))&=T(w_2\times 1+1\times w_2+\mu\times\mu)\\
&=T(w_2)\times 1+1\times T(w_2)+U(\mu)\times T(\mu)+T(\mu)\times U(\mu)\\
&=(T(w_2)+l\,U(\mu))\times 1+1\times(T(w_2)+k\,U(\mu)),\\
(\scr{L}_{k,l}\Delta)^*(U(\mu))&=U(\Delta^*\mu)=U(\mu)\times 1+1\times U(\mu),
\end{aligned}
\end{equation*}
从而
\begin{equation*}
\begin{aligned}
&(\scr{L}_{k,l}\Delta)^*\bigl(T(w_2)+(k+l-1)\,U(\mu)\bigr)\\
&\qquad=(T(w_2)+(k-1)\,U(\mu))\times 1+1\times(T(w_2)+(l-1)\,U(\mu)).
\end{aligned}
\end{equation*}
鉴于\eqref{eq:free-loop-cohomology} 中\trGenerator{}的表达式，上面的计算还表明：只要
两个\trVectorSpace{} $V,W$ 中至少一个的维数 $\ge 6$，$\scr{L}_{k,l}\Delta$ 就在
$H^1(-;\ZZ/2)$ 上诱导一个\trInjective{}映射（事实上 $\ge 4$ 就够了，但我们不坚持
这一点）。
(iii) \emph{$n=1$ 的计算}。设 $n=1$，此时\trMaslov{}数 $\mu(\rho)$ 化归为
\cref{sub:n-operators-line-bundles} 中考虑的卷绕数。例如，若 $S=D$ 且 $\mu(\rho)=-1$，
则由\cref{thm:injective-negative-degree} 知 $D_{S,\rho}$ 是\trInjective{}的，进而由%
\trIndex{}公式它是\trInvertible{}的。接下来考虑 $S=D$、$\mu(\rho)=0$ 的情形。
\eqref{eq:winding-number-count} 表明 $\ker(D_{S,\rho})$ 中任何 $X$ 都是一个处处非零的
函数。这意味着任意两个\trKernel{}元素都成比例，因为它们的商将是 $S$ 上一个取实边界
值的\trHolomorphic{}函数，从而必为常数。这意味着在边界点上的取值诱导一个%
\trCanonical{}\trIsomorphism{} $\ker(D_{S,\rho})\cong\rho(0)$。另一方面，\trIndex{}公式表明
该算子是满的，故同样有 $\det(D_{S,\rho})\cong\rho(0)$。
(iv) \emph{一个高维情形}。考虑 $\scr{L}_0Gr(V)$，$n\ge 3$，并设 $S=D$。
\eqref{eq:free-loop-cohomology} 中的第二个\trDirectSummand{}来自常值\trLoop{}子空间
$Gr(V)\subset\scr{L}_0Gr(V)$。若 $\rho$ 是常值的，则 $(E,F_\rho)$ 显然分解为一维%
\trDirectSummand，而 (iii) 中的论证随之表明
$\det(D_{S,\rho})\cong\lambda^{\mathrm{top}}(\rho(0))$，这意味着\eqref{eq:det-bundle-w1}
的两侧限制到 $Gr(V)$ 上时相合。另一方面，可以构造一个映射 $S^2\to Gr(V)$，它与 $w_2$
的\trPairing{}非平凡，而若把它看作一条\trLoop{} $S^1\to\Omega\,Gr(V)$，它就给出 $S^1$ 上
一个非平凡的\trDeterminantLine\trLineBundle。这在 [60, 第 6 章] 或 [34] 中有详细的展开，
此处不拟重复。两个观察合在一起表明：
\begin{equation*}
w_1(\underline{\det}_S)=T(w_2)+U(\mu)\quad\text{在 }\scr{L}_0Gr(V)\text{ 上}.
\end{equation*}
有了这些，\eqref{eq:det-bundle-w1} 可按如下思路证明。(i) 表明只需考虑 $S=D$。由 (iii)
已知该等式在 $\scr{L}_0Gr(\CC)$、$\scr{L}_{-1}Gr(\CC)$ 上成立。鉴于 (ii)，若等式在
$\scr{L}_kGr(V\oplus\CC)$ 上成立，则它在 $\scr{L}_kGr(V)$ 与 $\scr{L}_{k+1}Gr(V)$ 上也
成立；此外，当 $\dim(V)\ge 6$ 时还可以反向论证。从 (iv) 出发并利用这些观察，一般情形
便随之立得。
\end{proof}
\begin{remark}\label{thm:kuiper-bott-remark}
把\trHilbertSpace{}族 $W^{1,2}(D,E,F_\rho)$ \trTrivialization{}之后——这经 Kuiper 定理
总是可能的——诸 $D_{D,\rho}$（保持 $\rho(0)$ 不动）便定义了一个映射
$\Omega(\mathrm{U}_n/\mathrm{O}_n)\to\scr{F}(H,H)$。如 [34] 所解释的，应当把它（在%
\trIndex{}可相差一个常数的意义下）看作对\trBott{}\trPeriodic{}性%
\trHomotopyEquivalence{} $\scr{F}(H,H)\cong\ZZ\times B\mathrm{O}_\infty\cong
\Omega(\mathrm{U}_\infty/\mathrm{O}_\infty)$ 的一个近似。
\end{remark}
